% ============================================================
%  part1/introduction.tex  —  Introduction to Part I
% ============================================================

\chapter*{Introduction to Part I}
\addcontentsline{toc}{chapter}{Introduction to Part I}
\markboth{Introduction to Part I}{Introduction to Part I}

Part~I introduces the central thesis of this book through seven
examples that require no prior mathematics beyond basic algebra.
The examples are chosen because every reader already knows,
intuitively, that observation is not the same as reality.
The purpose of Part~I is to make that intuition precise.

\begin{tcolorbox}[
  title={\textbf{Established So Far}},
  colback=boxgreen,colframe=green!50!black,
  fonttitle=\small\sffamily
]
Nothing has been assumed beyond familiarity with the idea that
representations simplify.
No prior mathematics is required to enter Part~I.
\end{tcolorbox}

\bigskip

\begin{tcolorbox}[
  title={\textbf{New Ideas Introduced in This Part}},
  colback=boxblue,colframe=blue!40!black,
  fonttitle=\small\sffamily
]
\begin{itemize}[leftmargin=1.5em,itemsep=2pt]
  \item \textbf{Projection} $\pi : X \to M$: any map that
    preserves some structure while discarding the rest.
  \item \textbf{Kernel} $\ker(\pi)$: the pairs a projection
    fails to distinguish; the measure of what is lost.
  \item \textbf{Quotient space} $X/{\sim_\pi}$: the territory
    as the map represents it.
  \item \textbf{Admissibility class} $\pi^{-1}(m)$: the set of
    originals consistent with an observation.
  \item \textbf{Reconstruction error} $E(x) = d(x, G(F(x)))$:
    the fidelity of reconstruction without exactness.
  \item \textbf{Bidirectional information loss}: both
    coarse-graining and over-resolution destroy information.
  \item \textbf{Admissibility as task-relative}: the right
    observation is not the most detailed one.
  \item \textbf{Detector equivalence}: physically distinct
    states that produce identical measurements.
  \item \textbf{The Ocularum}: the structured observational
    domain of an observer.
  \item \textbf{Horizon-relative reconstruction}: different
    observers reconstruct different models of the same whole.
\end{itemize}
\end{tcolorbox}

\bigskip

\begin{tcolorbox}[
  title={\textbf{Dependencies}},
  colback=boxgray,colframe=gray!50!black,
  fonttitle=\small\sffamily
]
Part~I has no mathematical prerequisites.
Readers familiar with basic set theory will find the notation
straightforward.
All formal definitions are introduced from scratch.
\end{tcolorbox}

\bigskip

\begin{tcolorbox}[
  title={\textbf{What Remains Open After This Part}},
  colback=boxorange,colframe=orange!60!black,
  fonttitle=\small\sffamily
]
Part~I establishes the \emph{conceptual} framework.
It does not yet provide the mathematical machinery needed
to make admissibility precise, compute obstruction classes,
or derive the RSVP dynamics.
Those tools come in Parts~II--V\@.

In particular, the following questions raised in Part~I
are not answered until later:
\begin{itemize}[leftmargin=1.5em,itemsep=2pt]
  \item When can multiple local observations be assembled into
    a global reconstruction? (Chapter~\ref{ch:sheaf})
  \item What mathematical object measures the failure of
    global reconstruction? (Chapter~\ref{ch:monoturn})
  \item How does the admissibility of an observation depend
    on the governing dynamics? (Part~IV)
  \item What is the precise relationship between the RSVP
    fields and the admissibility framework? (Part~V)
\end{itemize}
\end{tcolorbox}

\bigskip

\noindent
The single question that Part~I raises and the rest of the book
answers:

\principle{
  Given a projection with nontrivial kernel,\\[4pt]
  what global structure can be reconstructed from what remains?
}
